\chapter*{Введение}							% Заголовок
\addcontentsline{toc}{chapter}{Введение}	% Добавляем его в оглавление
\label{intro}

Модель смеси экспертов (Mixture-of-Experts, MoE) --- это композитная
модель, в которой вход распределяется механизмом маршрутизации
(gating) между несколькими специализированными подмоделями-экспертами,
а итоговый ответ формируется как взвешенная комбинация их выходов
\cite{jacobs1991adaptive, jordan1994hierarchical, jacobs1997bayesian}.
Случаи, когда эксперты являются классическими статистическими моделями,
хорошо изучены в литературе \cite{piwko2025divide}, в частности, когда
эксперты заданы логистическими регрессиями \cite{nguyen2023general} или
гауссовскими моделями \cite{nguyen2023demystifying, rasmussen2001infinite}.
В последние годы, с ростом интереса к нейронным сетям, конструкция MoE
была перенесена на архитектуры глубокого обучения \cite{eigen2013learning}.
Высокая вычислительная эффективность достигается за счёт разреженной
маршрутизации, когда вход направляется лишь к небольшому фиксированному
числу экспертов \cite{shazeer2017outrageously, riquelme2021scaling, fedus2022switch};
особенно широкое распространение MoE получили в трансформерных моделях,
где размер модели играет ключевую роль
\cite{jiang2024mixtral, qiu2025demons, dai2024deepseekmoe, mu2025comprehensive}.

Параллельно методы поиска нейросетевой архитектуры (Neural Architecture
Search, NAS) развились в самостоятельную, хорошо изученную область
\cite{ren2021comprehensive, nayman2019xnas, zoph2018learning}. Цель NAS ---
автоматический выбор оптимальной архитектуры из широкого множества
кандидатов, что снимает необходимость в ручном подборе и экспертных
знаниях \cite{liu2018darts}. Сложились две основные парадигмы. Первая ---
вычислительно затратный многократный (many-shot) поиск с несколькими
циклами обучения и последующим выбором кандидата методами обучения с
подкреплением \cite{zoph2016neural, baker2016designing} или
эволюционными алгоритмами \cite{real2019regularized, yang2020cars}.
Вторая --- одношаговый (one-shot) поиск, где обучается единая
суперсеть, а затем из неё выбирается оптимальная подсеть
\cite{liu2018darts, chen2020stabilizing, you2020greedynas}.

Несмотря на успехи обеих методологий по отдельности, применение методов
NAS именно к архитектурам MoE остаётся слабо исследованной областью. В
отличие от классического подхода, в котором все эксперты имеют
\emph{одинаковую} архитектуру \cite{fedus2022switch}, в настоящей работе
допускается \emph{архитектурная гетерогенность} экспертов. Мы полагаем,
что такая гибкость повышает эффективность итоговой модели, позволяя
каждому эксперту специализироваться под свой тип данных. Близкие по
постановке работы немногочисленны: в \cite{jawahar2022automoe} быстрая
MoE-архитектура ищется эволюционно с наследованием весов из суперсети; в
\cite{han2024cmn} строится фреймворк для эффективного инференса MoE под
конкретное оборудование; наиболее близкая к нашей работа
\cite{mecharbat2025moenas} варьирует число экспертов в слоях. В отличие
от неё, мы оптимизируем сами архитектуры экспертов, а не только их
количество.

Архитектура MoE по своей природе хорошо подходит для данных с выраженной
кластерной структурой. Существуют теоретические работы, показывающие на
относительно простых примерах справедливость этого утверждения
\cite{chen2022towards}, а практические исследования демонстрируют
эффективность MoE на мультимодальных датасетах
\cite{chai2022mixture, guo2018multi, hu2025multi}. Тем не менее,
механизмы, благодаря которым эксперты выравниваются с различными
модальностями данных, изучены недостаточно.

В связи с этим формулируется следующий исследовательский вопрос:
\textit{как организовать условия, при которых архитектуры экспертов
адаптируются под конкретные кластеры данных?} В настоящей работе
предлагается фреймворк структурно-зависимого поиска архитектуры
(Structure-Aware NAS, SA-NAS), который совместно определяет архитектуры
экспертов и разбиение кластеров данных между ними. Ключевыми элементами
метода являются суррогатная модель, предсказывающая качество архитектуры
на заданном подмножестве данных без её фактического обучения, и
алгоритм оптимизации правдоподобия смеси на основе обобщённого
EM-алгоритма с адаптивно уточняемым суррогатом (Surrogate-assisted
Generalized EM, SGEM).

\medskip
\textbf{Основные результаты работы.}
\begin{itemize}
    \item \textbf{Постановка задачи.} Предложена новая постановка задачи
    поиска архитектуры для модели смеси экспертов как задачи максимизации
    кластерно-зависимого правдоподобия; дано её теоретическое
    обоснование --- показана эквивалентность оценке максимального
    правдоподобия по неполным данным для латентной модели смеси.
    \item \textbf{Алгоритм.} Предложен алгоритм решения поставленной
    задачи, основанный на EM-алгоритме с адаптивно уточняемой
    суррогатной моделью.
    \item \textbf{Гарантия сходимости.} Для предложенного алгоритма
    доказана теорема о сходимости (Теорема~\ref{thm:sgem}).
    \item \textbf{Датасеты и эксперименты.} Построены датасеты, на
    которых у подходов на основе смеси экспертов есть явное преимущество,
    и на них проведены сравнительные эксперименты с базовыми методами.
\end{itemize}

\medskip
\textbf{Обозначения.}\quad Жирными строчными буквами $\mathbf{y}$
обозначаются векторы, жирными прописными $\mathbf{X}$ --- матрицы.
Рассматривается задача обучения с учителем: задан датасет
$\mathcal{D} = (\mathbf{X}, \mathbf{y})$, где $\mathbf{X}$ ---
матрица признаков, $\mathbf{y}$ --- вектор целевых переменных. Данные
разбиты на обучающую и валидационную подвыборки, для которых заданы
функции потерь $\mathcal{L}_{\text{train}}$ и
$\mathcal{L}_{\text{val}}$ соответственно. Индекс
$n \in \{1, \ldots, N\}$ нумерует объекты, $m \in \{1, \ldots, M\}$ ---
кластеры данных, $k \in \{1, \ldots, K\}$ --- экспертов.
